\section{Introduction}
\label{sec:introduction}

Crude oil remains the most systemically consequential commodity in the
global economy. Despite accelerating investments in renewable capacity,
petroleum and its derivatives still account for roughly 30 percent of
primary energy consumption worldwide \citep{iea2025} and anchor the
fiscal architecture of more than two dozen sovereign states, from the
Gulf monarchies to Nigeria, from Russia to Venezuela. The transmission
mechanism is direct: supply disruptions propagate into headline
inflation, force central bank recalibrations, widen sovereign credit
spreads, and spill over into equity volatility and foreign-exchange
markets \citep{kilian2016,barskykilian2004}. The strategic interactions
among the world's three dominant production blocs: the Organization of
the Petroleum Exporting Countries (OPEC), US shale producers, and the
Russian Federation determine not merely the barrel price on trading
screens, but the macroeconomic trajectory of importing and exporting
nations alike. Understanding these interactions is a question at the
nexus of macroeconomics, geopolitics, and financial stability.

The urgency of that question has sharpened dramatically since late
February 2026. On 28 February, the United States and Israel launched
strikes against Iran, precipitating the closure of the Strait of Hormuz,
a 21-mile-wide chokepoint through which approximately 20 percent of
globally traded petroleum transits under normal conditions
\citep{crs2026}. The Iranian Revolutionary Guard declared the strait
closed, attacked merchant vessels, and laid sea mines, reducing shipping
traffic by more than 90 percent from pre-conflict levels. The EIA reports
that crude and fuel flows through the strait declined by nearly
6 million barrels per day in the first quarter of the conflict
\citep{eia2026}. Brent crude, which traded near \$65/bbl as recently as
November 2025, surged past \$116/bbl in early May 2026 before settling
around \$106--108/bbl amid uncertain ceasefire negotiations. Saudi
Aramco's CEO warned on the company's first-quarter earnings call that
the market is losing roughly 100 million barrels of supply per week and
that normalisation could extend into 2027 if the strait remains blocked
past mid-June \citep{cnbc2026}. The International Energy Agency (IEA)
has cautioned that global observed inventories fell at a record pace of
approximately 4 million barrels per day in March and April, and that the
market could remain severely undersupplied until October even if the
conflict ends next month \citep{iea2026omr}.

Superimposed on this acute supply shock is a structural fracture in
cartel governance. The United Arab Emirates (UAE) notified OPEC of its
decision to withdraw from OPEC and OPEC+ on 28 April 2026, with the exit
taking effect on 1 May 2026, after fifty-nine years of membership, the
most significant departure in the organisation's history
\citep{emirates2026}. The UAE's decision followed years of
friction over quota allocations that held its production roughly
30 percent below its 4.85 million barrel-per-day capacity
\citep{woodmac2026}. Abu Dhabi's exit follows Angola's departure in 2024
and Qatar's in 2019, establishing a pattern of defection that game theory
predicts with precision: when the temptation payoff from unilateral
expansion exceeds the discounted value of continued cooperation, the
cartel unravels \citep{fudenbergmaskin1986}. On the opposing flank of
supply, the US shale sector is approaching a structural plateau. The
EIA's \emph{Annual Energy Outlook 2026} projects that US crude
production will decline from its 2025 record of 13.6 million barrels per
day, as mature plays exhibit declining well productivity and capital
discipline continues to prioritize shareholder returns over volume
growth \citep{eia2026aep,dallasfed2026}. Russia, meanwhile, occupies a
paradoxical position: Western sanctions constrain its investment and
logistics, yet the Hormuz-induced price spike elevates the very revenue
those sanctions were designed to suppress \citep{ft2026}.

This paper investigates these strategic interactions through a
progressive computational framework that integrates classical game
theory with modern reinforcement learning. We pursue four interlocking
questions. First, what are the equilibrium outcomes of strategic
quantity competition among asymmetric oil producers under alternative
market structures? Second, can tacit collusion be sustained, and under
what conditions of patience, monitoring, and demand uncertainty does it
survive or collapse? Third, and this is the central empirical
contribution, does tacit collusion emerge spontaneously when autonomous
agents learn production strategies through trial and error, without
programmed coordination, without knowledge of the demand model, and
observing only the market price? Fourth, which production strategies
survive evolutionary selection in a population of heterogeneous
producers?

We construct an original Python codebase calibrated to a stylized
oil-market model with linear inverse demand $P(Q) = 140 - Q$ and
asymmetric marginal costs reflecting the structural heterogeneity among
US shale (\$45/bbl, consistent with the Dallas Fed Energy Survey
Q1 2026, \citealp{dallasfed2026}); OPEC (\$20/bbl, reflecting Gulf
weighted-average lifting costs per BP Statistical Review,
\citealp{bp2022}); and Russia (\$35/bbl, per IMF Article~IV,
\citealp{imf2021}).

The analytical pipeline escalates through seven conceptual layers,
detailed as a sixteen-step reading guide in
\appref{app:reading-guide}: (i)~static
Cournot competition establishing the Nash benchmark; (ii)~Stackelberg
leadership quantifying OPEC's first-mover advantage; (iii)~repeated
games and the Folk Theorem formalizing cartel sustainability;
(iv)~Shapley values addressing coalition stability;
(v)~\citet{greenporter1984} imperfect monitoring under stochastic
demand; (vi)~single-agent and multi-agent Q-learning under price-only
observation; and (vii)~evolutionary dynamics via replicator equations.
Each layer addresses a limitation of the preceding one, and all share a
common quantitative metric, which is the collusion index, defined as
\begin{equation}
\mathrm{CI} \;=\; \frac{P^{\,\text{Converged}} - P^{\,\text{Nash}}}{P^{\,\text{Cartel}} - P^{\,\text{Nash}}},
\label{eq:ci-intro}
\end{equation}
that enables direct cross-framework comparison.

The contribution is threefold. Methodologically, we are, to our
knowledge, the first to embed multi-agent Q-learning within the
Green-Porter imperfect monitoring framework and apply it to a
calibrated energy-market oligopoly, a design that extends the
Bertrand-focused experiments of \citet{calvano2020} to Cournot quantity
competition with asymmetric costs. Empirically, we provide computational
evidence that tacit collusion is an emergent property of autonomous
production-setting agents, with the collusion index varying
systematically with the discount factor, exploration rate, and learning
rate in directions consistent with the Folk Theorem's predictions. For
policy, these findings speak directly to the accelerating regulatory
debate around algorithmic collusion: California's AB 325, effective
January 2026, targets algorithms that facilitate price alignment
\citep{dlapiper2025}; the European Commission has confirmed multiple
active investigations into algorithmic pricing \citep{nortonrose2025,loyensloeff2026};
and the US Department of Justice has warned of the risk of ``fully
automated cartels'' \citep{kavanagh2025}.
